\documentclass[tikz,border=1.5mm]{standalone}
\input{src/_preamble.tex}
\begin{document}
\begin{tikzpicture}[x=1mm,y=1mm]
  \tikzset{col/.style={draw=ink, fill=pale, rounded corners=2pt, minimum width=20mm, minimum height=7mm, font=\scriptsize\bfseries},
           role/.style={draw=ink, fill=white, rounded corners=1pt, text width=18mm, align=center, font=\tiny, inner sep=1.2mm, minimum height=7.5mm},
           from/.style={text width=20mm, align=center, font=\tiny, text=mid}}
  \def\cols{本体/0, 小脑/22, 感知/44, 大脑/66, 软件/88, 交付/110}
  \foreach \n/\x in \cols \node[col] at (\x,40) {\n};
  % 本体
  \node[role] at (0,30) {结构工程师（4）};
  \node[role] at (0,21) {关节模组（5、6）};
  \node[role] at (0,12) {电机控制（5）};
  \node[role] at (0,3) {硬件与电源（7）};
  \node[role] at (0,-6) {嵌入式（7、9）};
  % 小脑
  \node[role] at (22,30) {运动控制·模型派（10、11）};
  \node[role] at (22,21) {强化学习·学习派（12、13）};
  \node[role] at (22,12) {仿真工程师（12）};
  % 感知
  \node[role] at (44,30) {感知（8）};
  \node[role] at (44,21) {SLAM 与定位（14）};
  \node[role] at (44,12) {导航（14）};
  % 大脑
  \node[role] at (66,30) {具身大模型（15）};
  \node[role] at (66,21) {模仿学习（15）};
  \node[role] at (66,12) {数据采集与遥操作（15）};
  % 软件
  \node[role] at (88,30) {机器人软件与中间件（9）};
  \node[role] at (88,21) {工具链与数据平台};
  % 交付
  \node[role] at (110,30) {测试与可靠性（4、17）};
  \node[role] at (110,21) {制造、工艺、供应链（16）};
  \node[role] at (110,12) {产品、售前、技术支持};
  \node[role] at (110,3) {安全工程师（17）};
  % 来源背景
  \draw[mid, line width=0.4pt] (-11,-14) -- (121,-14);
  \node[font=\tiny\bfseries, text=ink, anchor=west] at (-11,-17.5) {最顺的来源背景};
  \node[from] at (0,-24) {手机、家电、汽车结构；新能源三电；嵌入式};
  \node[from] at (22,-24) {自动驾驶规控；强化学习算法};
  \node[from] at (44,-24) {自动驾驶感知与定位；扫地机、AR};
  \node[from] at (66,-24) {大模型、多模态、视觉算法};
  \node[from] at (88,-24) {互联网后端；自动驾驶中间件};
  \node[from] at (110,-24) {工业自动化；消费电子与汽车制造；网络安全};
  \node[font=\tiny, text=mid, anchor=west] at (-11,-33) {括号里是本书对应的章};
\end{tikzpicture}
\end{document}
